\section{PCA and explained variance}
\label{sec:ex8}
\begin{problem}
The covariance matrix of centered data has eigenvalues $\lambda_1=6$, $\lambda_2=2$, $\lambda_3=1$, and $\lambda_4=1$.
\begin{enumerate}[label=\textbf{(\alph*)}]
\item Calculate the explained-variance percentage for each principal component.
\item Find the fewest components needed to preserve at least $85\%$ of the total variance.
\item Explain why principal directions are orthogonal and what the first component optimizes.
\item Explain why low reconstruction error under PCA need not improve classification.
\end{enumerate}
\end{problem}


\begin{concept}
For a covariance matrix $\Sigma$, each principal direction $v_j$ satisfies $\Sigma v_j=\lambda_j v_j$. Its eigenvalue is the variance of data projected onto that unit-length direction. The total variance is $\operatorname{tr}(\Sigma)=\sum_j\lambda_j$.
\end{concept}

\subq{(a) Explain the variance percentage of each component}
First calculate total variance:
\[
\lambda_{\rm total}=6+2+1+1=10.
\]
Then compute $100\lambda_j/10$:
\begin{center}\begin{tabular}{@{}rrrrr@{}}\toprule Component&PC1&PC2&PC3&PC4\\\midrule Eigenvalue&6&2&1&1\\ Explained variance&$60\%$&$20\%$&$10\%$&$10\%$\\ Cumulative&$60\%$&$80\%$&$90\%$&$100\%$\\\bottomrule\end{tabular}\end{center}

\subq{(b) Minimum number needed to retain at least $85\%$}
Find the smallest $q$ such that $\sum_{j=1}^{q}\lambda_j/\sum_{j=1}^{4}\lambda_j\geq0.85$:
\[
q=2:\ \frac{6+2}{10}=0.80<0.85,
\qquad q=3:\ \frac{6+2+1}{10}=0.90\geq0.85.
\]
Therefore $\boxed{q_{\min}=3}$. Two dimensions retain only $80\%$; three retain $90\%$.

\subq{(c) Why are principal directions orthogonal? What does PC1 optimize?}
A real covariance matrix is symmetric ($\Sigma^\top=\Sigma$) and positive semidefinite. The spectral theorem guarantees an orthonormal eigenbasis. To see why eigenvectors for \emph{distinct} eigenvalues are orthogonal, let $\Sigma v_i=\lambda_i v_i$ and $\Sigma v_j=\lambda_j v_j$; then
\[
\lambda_j v_i^\top v_j=v_i^\top\Sigma v_j
=(\Sigma v_i)^\top v_j=\lambda_i v_i^\top v_j,
\]
so $(\lambda_i-\lambda_j)v_i^\top v_j=0$. If $\lambda_i\neq\lambda_j$, it follows that $v_i^\top v_j=0$. For repeated eigenvalues, such as $\lambda_3=\lambda_4=1$, choose an orthonormal basis within that eigenspace.

PC1 maximizes the variance of the projected centered data among all unit-length directions:
\[
v_1\in\argmax_{\norm{v}_2=1}\Var(Xv)
=\argmax_{\norm{v}_2=1}v^\top\Sigma v.
\]
The Rayleigh quotient achieves its maximum $\lambda_1=6$ at a leading eigenvector.

\subq{(d) Why might excellent PCA reconstruction fail to help classification?}
PCA is \emph{unsupervised}: it preserves directions with high overall variance, not necessarily those that separate labels. If discriminative information lies in a low-variance direction, dimensionality reduction can remove it while retaining most reconstruction variance. Classification quality must therefore be checked using a held-out or cross-validated supervised metric, with PCA fitted within training folds.
\begin{answer}
Explained fractions $=60\%,20\%,10\%,10\%$. Preserve at least $85\%$ with $\boxed{3}$ principal components.
\end{answer}
